\scanpage{30}
% archive sheet number: 29
% source top heading
\uline{$n$-Quasi-polyèdres réguliers en coordonnées can.}

% source lines 1-2
\[
\left\{\begin{aligned}
&\sigma_0,\ldots,\sigma_n\in\mathrm{Aff}(E)\qquad E\ \text{affine de rang }n+1\\
&\sigma_j\sigma_i=\sigma_i\sigma_j\quad\text{si }0\leq i\leq i+2\leq j\leq n
\end{aligned}\right.\tag*{(1)}
\]
% source line 3
\[
E_0=\{s_0\}\subset E_1\subset\cdots\subset E_n\qquad\text{drapeau de }E\tag*{(\unsure{1|2})}
\]
% source lines 4-5
\[
\left\{\begin{aligned}
\sigma_iE_j&=E_j\quad\text{si }j\ne i\\
\sigma_iE_i&\ne E_i
\end{aligned}\right.\tag*{(3)}
\]
% source lines 6-10
\[
\left\{\begin{aligned}
&s_0\\
&s_1=\sigma_0s_0\\
&\cdots\\
&s_i=\sigma_{i-1}s_{i-1}\qquad(1\leq i\leq n+1)\\
&\cdots\\
&s_{n+1}=\sigma_ns_n
\end{aligned}\right.\tag*{(4)}
\]
% source lines 11-14
\uline{Prop} $(s_0,\ldots,s_{n+1})$ est une base affine de $E$, et
\[
E_i=\unsure{\mathrm{Env.aff}|\mathrm{Ens.aff}}(s_0,\ldots,s_i)\tag*{(5)}
\]
\[
\sigma_ks_i=s_i\quad\text{si }0\leq i<k\leq n\deleted{\unsure{+1}}\tag*{(6)}
\]
\inserted{\uline{pour $\unsure{i}$, $0\leq i\leq n+1$}}

% source lines 15-17
Par induction \unsure{(on va prouver} que $(s_0,\ldots,s_i)$ est aff\textsuperscript{t} libre et (5), (6). Trivial si $i=0$. Supposons \inserted{$i\geq1$ et} prouvé pour $j<i$. Il faut prouver
% source lines 18-20
\[
\left\{\begin{aligned}
&s_i\notin E_{i-1}(=\mathrm{Env}(s_0,\ldots,s_{i-1})\ \text{par hyp. réc})\\
&E_i=\mathrm{Env}(E_{i-1},s_i)
\end{aligned}\right.
\]
\inserted{$\sigma_{\unsure{h|k}}s_i=s_i$ pr $i<h\leq n$}
% source lines 21-23
Si on avait $s_i\in E_{i-1}$, i.e. $\sigma_{i-1}s_{i-1}\in E_{i-1}$, comme $\sigma_{i-1}s_j=s_j$ pour $0\leq j\leq i-2$ et $E_{i-1}=\mathrm{Env}(s_0,\ldots,s_{i-1})$ par réc, on aurait
\[
\sigma_{i-1}E_{i-1}=\mathrm{Env}(\underset{\substack{\Vert\\s_0}}{\sigma_{i-1}s_0},\ldots,\underset{\substack{\Vert\\s_{i-2}}}{\sigma_{i-1}s_{i-2}},\underset{\substack{\uparrow\\E_{i-1}}}{\sigma_{i-1}s_{i-1}})\subset E_{i-1}
\]
% source lines 25-28
donc $\sigma_{i-1}E_{i-1}=E_{i-1}$, contrairement à (3 b), OK.

On a, comme $\sigma_{i-1}E_i=E_i$ et $s_{i-1}\in E_{i-1}$ (réc) $\subset E_i$, $\sigma_{i-1}(s_{i-1})\subset E_i$ i.e. $s_i\in E_i$. Comme $s_i\notin E_{i-1}$, on a donc $E_i=\mathrm{Env}(E_{i-1},s_i)$. Enfin, on a pour \deleted{$h>i$}\inserted{$i<h\leq n$}

\[
\sigma_hs_i=\sigma_h\sigma_{i-1}s_{i-1}=\sigma_{i-1}(\sigma_hs_{i-1})\overset{\text{réc}}{=}\sigma_{i-1}(s_{i-1})=s_i
\]
OK.
